Coincidences among small mesh patterns have been considered by
\textcite{DBLP:journals/corr/ClaessonTU14}, in which the authors use the Simultaneous Shading Lemma,
a closure result and one worked out special case to fully classify coincidences among mesh patterns
of length 2.

Recall that two patterns \(\lambda\) and \(\gamma\) are considered \emph{coincident} if the set of permutations
that avoid \(\lambda\) is the same as the set of permutations that avoid \(\gamma\),
\ie \(\av{\lambda} = \av{\gamma}\). Equivalently we can say that they have the same set of
\emph{containers}, \ie \(\cont{\lambda} = \cont{\gamma}\).

We will consider the avoidance sets \(\av{\{\pi,p\}}\) where \(\pi\) is a classical pattern of length 3
and \(p\) is a mesh pattern of length 2 to establish sufficient conditions for two such
sets to be coincident. The classical pattern \(\pi\) will be fixed and called
the \emph{dominating pattern}.

In order to describe the rules it is useful to have a notion for inserting points, ascents, and
descents into a mesh pattern.
\begin{definition}
\label{def:ap}
Let \(p=(\tau,R)\) be a mesh pattern of length \(n\) such that \(\boks{i,j}\notin R\). We define
a mesh pattern \(p^\boks{i,j} = (\tau^\prime,R^\prime)\) of length \(n+1\) as the pattern where a
point is \emph{inserted} into the box \(\boks{i,j}\) in \(G(p)\). Formally the new underlying
classical pattern is defined by
\begin{equation*}
\tau^\prime(k) = \begin{cases}
    j+1 & \text{if } k = i+1\\
    \tau(k) & \text{if } \tau(k)\le j \text{ and }k\le i\\
    \tau(k)+1 & \text{if } \tau(k)> j \text{ and }k\le i\\
    \tau(k-1) & \text{if } \tau(k)\le j \text{ and }k> i+1\\
    \tau(k-1)+1 & \text{if } \tau(k)> j \text{ and }k> i+1\\
\end{cases}
\end{equation*}
While the mesh becomes
\begin{equation*}
\begin{aligned}
R^\prime=&\{\boks{k,\ell}\mid k\le i, \ell\le j, \boks{k,\ell}\in R\} \cup\\
&\{\boks{k,\ell}\mid k\le i, \ell > j, \boks{k,\ell-1}\in R\}\cup\\
&\{\boks{k,\ell}\mid k > i, \ell \le j, \boks{k-1,\ell}\in R\}\cup\\
&\{\boks{k,\ell}\mid k > i, \ell > j, \boks{k-1,\ell-1}\in R\}\\
\end{aligned}
\end{equation*}
\end{definition}
In addition, we give the following definitions:

\begin{definition}
Let \(p=(\tau,R)\) be a mesh pattern of length \(n\) such that \(\boks{i,j}\notin R\) and
\(p^\boks{i,j} = (\tau^\prime,R^\prime)\) is as defined in \DefinitionRef{def:ap}. We define
the following four modifications of \(p^\boks{i,j}\).
\begin{align*}
p^\boksup{i,j} &= (\tau^\prime, R^\prime \cup \{\boks{i,j+1},\boks{i+1,j+1}\})\\
p^\boksright{i,j} &= (\tau^\prime, R^\prime \cup \{\boks{i+1,j},\boks{i+1,j+1}\})\\
p^\boksdown{i,j} &= (\tau^\prime, R^\prime \cup \{\boks{i,j},\boks{i+1,j}\})\\
p^\boksleft{i,j} &= (\tau^\prime, R^\prime \cup \{\boks{i,j},\boks{i,j+1}\})\\
\end{align*}
\end{definition}
Informally, these are considering the topmost, rightmost, leftmost, or bottommost point in
\(\boks{i,j}\). We collect the resulting mesh patterns in a set
\begin{equation*}
p^\boksall{i,j}=\left\{
p^\boks{i,j},
p^\boksright{i,j},
p^\boksleft{i,j},
p^\boksup{i,j},
p^\boksdown{i,j}\right\}
\end{equation*}
See \FigureRef{fig:addp} for an example of adding a point into a mesh pattern.

\begin{figure}
\begin{equation*}
p = \pattern{}{1,2}{0/1,2/2} \quad
p^\boks{2,1}= \pattern{}{1,3,2}{0/1,0/2,2/3,3/3}
\quad
p^\boksup{2,1}
= \pattern{}{1,3,2}{0/1,0/2,2/2,2/3,3/2,3/3}
\end{equation*}
\caption{The result of inserting a point into \(p=(12,\{(0,1),(2,2)\})\)}
\label{fig:addp}
\end{figure}

\begin{definition}
Let \(p=(\tau,R)\) be a mesh pattern of length \(n\) such that \(\boks{i,j}\notin R\). We define
a mesh pattern \(p^{\boks{i,j}_a} = (\tau^\prime,R^\prime)\) (\(p^{\boks{i,j}_d}\)) of length \(n+2\)
as the pattern where an ascent (descent) is \emph{inserted} into the box \(\boks{i,j}\) in \(G(p)\).
Formally the new underlying classical pattern is defined by
\begin{equation*}
\tau^\prime(k) = \begin{cases}
    j+t & \text{if } k = i+t,t\in\{1,2\}\\
    \tau(k) & \text{if } \tau(k)\le j \text{ and }k\le i\\
    \tau(k)+2 & \text{if } \tau(k)> j \text{ and }k\le i\\
    \tau(k-2) & \text{if } \tau(k)\le j \text{ and }k> i+2\\
    \tau(k-2)+2 & \text{if } \tau(k)> j \text{ and }k> i+2\\
\end{cases}
\end{equation*}
The ordering of the top branch determines whether an ascent(or descent) is added.
The mesh becomes
\begin{equation*}
\begin{aligned}
R^\prime=&\{\boks{k,\ell}\mid k\le i, \ell\le j, \boks{k,\ell}\in R\} \cup\\
&\{\boks{k,\ell}\mid k\le i, \ell > j, \boks{k,\ell-2}\in R\}\cup\\
&\{\boks{k,\ell}\mid k > i, \ell \le j, \boks{k-2,\ell}\in R\}\cup\\
&\{\boks{k,\ell}\mid k > i, \ell > j, \boks{k-2,\ell-2}\in R\}\cup\\
&\{\boks{i+1,j},\boks{i+1,j+1},\boks{i+1,j+2}\}
\end{aligned}
\end{equation*}
\end{definition}
An example of adding an ascent to a mesh pattern can be seen in \FigureRef{fig:adda}.
\begin{figure}
\begin{equation*}
p = \pattern{}{2,3,1}{0/0,1/0,1/3,
                      3/0,3/1,3/3} \quad
p^{\boks{1,1}_a}= \pattern{}{4,2,3,5,1}{0/0,
                        1/0,1/5,
                        2/0,2/1,2/2,2/3,2/5,
                        3/0,3/5,
                        5/0,5/1,5/2,5/3,5/5}
\end{equation*}
\caption{The result of inserting an ascent into \(p=(231,\{(0,0),(1,0),(1,3),(3,0),(3,1),(3,3)\})\)}
\label{fig:adda}
\end{figure}

We now attempt to fully classify coincidences in families characterised by avoidance
of a classical pattern of length \(3\) and a mesh pattern of length \(2\), that is
finding and explaining all coincidences where \(\av{\{p,m\}} = \av{\{p,m^\prime\}}\).

It can be easily seen that in order to classify coincidences one need only
consider coincidences within the family of mesh patterns with the same underlying
classical pattern, this is due to the fact that \(\perm{2,1}\in\av{\mperm{1,2}{R}}\)
and \(\perm{1,2}\in\av{\{\mperm{2,1}{R}\}}\) for all mesh-sets \(R\).

We know that there are a total of \(512\) mesh-sets for each underlying classical
pattern. By use of the previous results of \textcite{DBLP:journals/corr/ClaessonTU14}
the number of coincidence classes can be reduced to \(220\).

\subsection{Coincidence classes of Av(\{321, (21, \textit{R})\})}
Through experimentation, considering avoidance of permutations of up to length \(11\), we discover
that there are at least \(29\) coincidence classes of mesh patterns with underlying classical
pattern \(21\).

\begin{proposition}[First Dominating Pattern Rule]
    \label{prop:dom1}
    Given two mesh patterns \(m_1 =(\sigma, R_1)\) and \(m_2 = (\sigma, R_2)\),
    and a dominating classical pattern \(\pi = (\pi,\emptyset)\) such that
    \(\setsize{\pi} \le \setsize{\sigma} + 1\), the sets \(\av{\{\pi,m_1\}}\) and
    \(\av{\{\pi,m_2\}}\) are coincident if

    \begin{enumerate}
        \item The symmetric difference, \(R_1 \triangle R_2\), is \(\{(a,b)\}\)
        \item \(\pi \preceq \sigma^\boks{a,b}\)\label{prop:dom1:cont}
    \end{enumerate}
\end{proposition}

This rule can be understood very in graphical form. In the
pattern in \FigureRef{fig:rule1} we can gain shading in two boxes since if
there is a point in any of these boxes we would get an occurrence of the
dominating pattern \(\perm{3,2,1}\).

\begin{figure}[htb]
    \begin{center}
    \begin{tikzpicture}[scale=\picscale]
        \modpattern{}{2,1}{1/0,1/1,1/2,2/2}
    \end{tikzpicture}
        \raisebox{2ex}{\(\mapsto\)}
    \begin{tikzpicture}[scale=\picscale]
        \filldraw (2.5,0.5) circle (3pt);
        \modpattern{}{2,1}{1/0,1/1,1/2,2/2,2/0}
    \end{tikzpicture}

        \caption{Visual depiction of first dominating pattern rule.}
        \label{fig:rule1}
    \end{center}
\end{figure}

In order to prove the proposition we must first make the following note.

\begin{note}
    \label{not:downcmesh}
    Let \(R_1 \subseteq R_2\). Then any occurrence of \((\tau, R_2)\) in a permutation
    is an occurrence of \((\tau,R_1)\).
\end{note}

\begin{proof}[Proof of \PropositionRef{prop:dom1}]
    We need to prove that \(\av{\{\pi,m_1\}} = \av{\{\pi,m_2\}}\).

    \noindent Assume without meaningful loss of generality that \(R_2 = R_1 \cup \{(a,b)\}\).
    Since \(R_1\) is a subset of \(R_2\), \NoteRef{not:downcmesh} states that
    \(\av{\{\pi,m_1\}} \subseteq \av{\{\pi,m_2\}}\)

    Now we consider a permutation \(\omega \in \av{\pi}\),
    containing an occurrence of \(m_1\). Consider placing a point in the
    region corresponding to the box \((a,b)\), regardless of where in this
    region we place the point by condition~\ref{prop:dom1:cont} of the
    proposition we create an occurrence of \(\pi\). Therefore there can be no
    points in this region, which implies this is an occurrence of \(m_2\). Hence every occurrence of
    \(m_1\) is in fact an occurrence of \(m_2\), and we have that
    \(\av{\{\pi,m_2\}} \subseteq \av{\{\pi,m_1\}}\).

    Taking both directions of the containment we can therefore draw the
    conclusion that \(\av{\{\pi,m_1\}} = \av{\{\pi,m_2\}}\).
\end{proof}

All coincidence classes of \(\av{\{\perm{3,2,1}, \mperm{2,1}{R}\}}\) can be
explained by application of \PropositionRef{prop:dom1}. Thus explaining the
29 experimental coincidence classes opserved.

\subsection{Equivalence classes of Av(\{231, (21, \textit{R})\}).}
By application of \PropositionRef{prop:dom1} we obtain \(43\) coincidence
classes. Experimentation shows that there are in fact at least \(39\) coincidence
classes, for example the following two patterns seem to be coincident in
\(\av{\perm{2,3,1}}\) but this is not explained by \PropositionRef{prop:dom1}.
\begin{equation*}
    m_1 = \pattern{}{2,1}{0/0,0/1,0/2,1/0,2/1} \text{ and } m_2 = \pattern{}{2,1}{0/0,0/1,0/2,1/0,1/1,2/1}
\end{equation*}

Consider an occurrence of \(m_1\) in a permutation in \(\av{\perm{2,3,1}}\), consisting
of elements \(x\) and \(y\). If the region corresponding
to the box \((1,1)\) is empty we have an occurrence of \(m_2\). Otherwise, if there is any
increase in this box then we would have an occurrence of \(\perm{2,3,1}\),
however, since we are in \(\av{\perm{2,3,1}}\) this is not possible. This box
must therefore contain a (non-empty) decreasing subsequence.
 This gives rise to the following lemma:

\begin{lemma}
    \label{lem:incdec}
    Let \(m =(\sigma, R)\) be a mesh pattern, where the box \((a,b)\) is not
    in \(R\), and \(\pi = (\pi,\emptyset)\) be a dominating classical pattern. If
    \(\pi \preceq p^{\boks{i,j}_a}\)\\(\(\pi \preceq p^{\boks{i,j}_d}\)),
    then in any occurrence of \(m\) in a permutation \(\varrho\), the region
    corresponding to the box \((a,b)\) can only contain an decreasing
    (increasing) subsequence of \(\varrho\).
\end{lemma}
The proof is analogous to the proof of \PropositionRef{prop:dom1}.

Going back to our example mesh patterns
\begin{equation*}
    \begin{tikzpicture}[scale=\picscale]
        \modpattern{}{2,1}{0/0,0/1,0/2,1/0,2/1}
        \filldraw (1.66,1.33) circle (3pt);
        \filldraw (1.33,1.66) circle (3pt);
    \end{tikzpicture}
\end{equation*}
we know that the region corresponding to the box \((1,1)\) contains a decreasing subsequence.
If we let \(z\) be the topmost point in this decreasing subsequence, then \(xz\)
is an occurrence of \(m_2\). This shows that our two example patterns are coincident.

This result generalises into the following rule for categorising coincidences
of mesh patterns in cases where there is a dominating classical pattern.

\begin{proposition}[Second Dominating Pattern Rule]
    \label{prop:dom2}
    Given two mesh patterns \(m_1 =(\sigma, R_1)\) and \(m_2 = (\sigma, R_2)\),
    and a dominating classical pattern \(\pi = (\pi,\emptyset)\) such that
    \(\setsize{\pi} \le \setsize{\sigma} + 2\), the sets \(\av{\{\pi,m_1\}}\) and
    \(\av{\{\pi,m_2\}}\) are coincident if

    \begin{enumerate}
        \item \(R_1 \triangle R_2 = \{(a,b)\}\)
        \item   \begin{enumerate}
                \item\label{prop:dom2:condc} \(\pi \preceq p^{\boks{i,j}_a}\) and
                        \begin{enumerate}
                            \item \((a+1,b) \in \sigma\) and \((a+1,b-1)\notin R\) and \\
                                \((x,b-1)\in R \implies (x,b) \in R \) (where \(x\neq a,a+1\)) and\\
                                  \((a+1,y)\in R \implies (a,y) \in R\) (where \(y\neq b-1,b\)).
                            \item \((a,b+1) \in \sigma\) and \((a-1,b+1)\notin R\) and \\
                                  \((x,b+1)\in R \implies (x,b) \in R\) (where \(x\neq a-1,a\)) and\\
                                  \((a-1,y)\in R \implies (a,y) \in R\) (where \(y\neq b,b+1\)).
                        \end{enumerate}
                    \item \(\pi \preceq p^{\boks{i,j}_d}\) and
                        \begin{enumerate}
                            \item \((a+1,b+1) \in \sigma\) and \((a+1,b+1)\notin R\) and \\
                                  \((x,b+1)\in R \implies (x,b) \in R\) (where \(x\neq a,a+1\)) and\\
                                  \((a+1,y)\in R \implies (a,y) \in R\) (where \(y\neq b,b+1\)).
                            \item \((a,b) \in \sigma\) and \((a-1,b-1)\notin R\) and \\
                                  \((x,b+1)\in R \implies (x,b) \in R\) (where \(x\neq a-1,a\)) and\\
                                  \((a-1,y)\in R \implies (a,y) \in R\)  (where \(y\neq b-1,b\)).
                        \end{enumerate}
                \end{enumerate}
    \end{enumerate}
\end{proposition}
\begin{proof}
    Assume without meaningful loss of generality that \(R_2 = R_1 \cup \{(a,b)\}\).
    By \NoteRef{not:downcmesh} we only need to show that an occurrence of \(m_1\)
    implies and occurrence of \(m_2\).
    We consider taking the first branch of every choice.
    Now consider a permutation \(\omega \in \av{\pi}\). Suppose
    \(\omega\) contains \(m_1\) and consider the region corresponding to
    \((a,b)\) in \(R_1\).If the region is empty, the occurrence of \(m_1\) is trivially an occurrence
    of \(m_2\).
    If the region is non-empty, then by \LemmaRef{lem:incdec} and
    condition~\ref{prop:dom2:condc} of the proposition it must contain
    a decreasing subsequence.
    We can choose the topmost point in the region to replace the corresponding
    point in the mesh pattern and the points from the subsequence are now in
    the box southeast of the point. The other conditions allow this to be done
    without points being present in regions that were shaded. Hence there are no
    points in the region corresponding to the box \((a,b)\) in the mesh
    pattern, and therefore we can shade this region. This implies that every
    occurrence of \(m_1\) in \(\av{\pi}\) is in fact an occurrence of \(m_2\)
    so \(\av{\{\pi,m_2\}} \subseteq \av{\{\pi,m_1\}}\).

    Similar arguments cover the remainder of the branches.
\end{proof}

This proposition essentially states that we slide all of the points in the box we desire
to shade diagonally, and chose the topmost/bottommost point to replace the original
point in the mesh pattern.

By applying both the First Dominating Pattern Rule and the Second Dominating
Pattern Rule fully explains coincidences of classes of the form
\(\av{\{\perm{2,3,1},\mperm{2,1}{R}\}}\), obtaining 39 coincidence classes of
mesh patterns, confirming experimental observations.

\subsection{Equivalence classes of Av(\{231, (12, \textit{R})\}).}
When considering the coincidence classes of \(\av{\{\perm{2,3,1},\mperm{1,2}{R}\}}\)
we first apply the two Dominating Pattern rules previously established.
Starting from \(220\) classes, application of the first Dominating Pattern rule
gives \(85\) classes. Following this with the second Dominating Pattern rule
reduces the number of classes to \(59\). However we know that there are patterns
where the coincidences are not explained by the rules given above.

For example the patterns
\begin{equation*}
    m_1 = \pattern{}{1,2}{0/2,2/0,2/1} \text{ and } m_2 = \pattern{}{1,2}{0/2,1/0,2/0,2/1}
\end{equation*}
are experimentally coincident.

Consider an occurrence of \(m_1\) in a permutation, if the region corresponding
to the box \((1,0)\) is empty then we have an occurrence of \(m_2\). Now look at
the case when this region is not empty, and consider choosing the rightmost
point in the region.

This gives us an occurrence of the following mesh pattern.
\begin{equation*}
    \pattern{}{2,1,3}{0/3,2/0,2/1,3/0,3/1,3/2}
\end{equation*}
By application of \PropositionRef{prop:dom1} we then achieve the following
mesh pattern
\begin{equation*}
    \pattern{2,3}{2,1,3}{0/3,1/3,2/0,2/1,3/0,3/1,3/2}
\end{equation*}
If we look at the highlighted points we see that the subpattern is an occurrence
of the mesh pattern that we originally desired. This gives rise to the following
rule:

\begin{proposition}[Third Dominating Pattern Rule]
    \label{prop:dom3}
    Given two mesh patterns \(m_1 =(\sigma, R_1)\) and \(m_2 = (\sigma, R_2)\),
    and a dominating classical pattern \(\pi = (\pi,\emptyset)\), the sets
    \(\av{\{\pi,m_1\}}\) and \(\av{\{\pi,m_2\}}\) are coincident if
    \begin{enumerate}
        \item \(R_1 \triangle R_2 = \{(a,b)\}\)
        \item\label{prop:dom3:condocc} One of the patterns in \((\sigma,R_1)^\boksall{a,b}\)
            is coincident with a mesh pattern containing an occurrence of
            \(\mperm{\sigma}{R_2}\) as a subpattern.
    \end{enumerate}
\end{proposition}
\begin{proof}
    Assume without meaningful loss of generality that \(R_2 = R_1 \cup \{(a,b)\}\).
    \NoteRef{not:downcmesh}, implies \(\av{\{\pi,m_1\}} \subseteq \av{\{\pi,m_2\}}\) as before.

    Now consider a permutation \(\sigma\) in \(\av{\pi}\) that contains an
    occurrence of \(m_1\). If the region corresponding to the box \((a,b)\)
    is empty then we have an occurrence of \(m_2\). If the region is non-empty
    then by condition~\ref{prop:dom3:condocc}  of the proposition there exists
    a pattern in \((\sigma,R_1)^\boksall{a,b}\) such that there exists an occurrence of a mesh pattern of
    length one longer than \(m_1\) in this position. This mesh pattern is
    coincident with another mesh pattern that contains an
    occurrence of \(m_2\). Hence, every occurrence of \(m_1\) leads to an occurrence
    of \(m_2\).  Thus \(\av{\{\pi,m_2\}} \subseteq \av{\{\pi,m_1\}}\) and the
    two patterns are coincident.
\end{proof}

By application of this rule we can reduce the number of classes in
\(\av{\{\perm{2,3,1},(\perm{1,2}, R)\}}\) to \(56\), which coincides with the
56 classes observed through experimentation.

\subsection{Equivalence classes of Av(\{321, (12, \textit{R})\}).}
When considering coincidences of mesh patterns with underlying classical
pattern \(\perm{1,2}\) in \(\av{\perm{3,2,1}}\) application of the previously
established rules give no coincidences. Through experimentation we discover
that there are \(7\) non-trivial coincidence classes (all others are singletons)
which can be explained through the
use of two different lines of reasoning. Since the number of coincidences is
so small we will reason for these coincidences without attempting to generalise
into concrete rules.

Intuitively it is easy to see why our previous rules have no
power here. There is nowhere that it is possible to add a single point to
gain an occurrence of \(\pi = \perm{3,2,1}\). It is also impossible to have
a position where addition of an increase, or decrease, provides extra
shading power.

The patterns
\begin{equation*}
    m_1 = \pattern{}{1,2}{0/0, 0/2, 1/1, 1/2, 2/1} \text { and } m_2 \pattern{}{1,2}{0/0, 0/1, 0/2, 1/1, 1/2, 2/1}
\end{equation*}
are equivalent in \(\av{\perm{3,2,1}}\). (There are 3 symmetries of these
patterns that are also equivalent to each other by the same reasoning.)

Consider the region corresponding to the box \((0,1)\) in any occurrence of \(m_1\),
in a permutation. By \LemmaRef{lem:incdec} it must contain an increasing subsequence.
If the region is empty then we have an occurrence of \(m_2\).
If there is only one point in the region we can choose this to replace the \(1\) in
the mesh pattern to get the required shading. If there is more than one point then
choosing the two leftmost points gives us the following mesh pattern.
\begin{equation*}
    \pattern{3,4}{2,3,1,4}{0/0, 0/1, 0/2, 0/3, 0/4,
                            1/0, 1/1, 1/2, 1/3, 1/4,
                            2/0, 2/1,2/2, 2/4,
                            3/1, 3/2, 3/3, 3/4,
                            4/1, 4/2, 4/3}
\end{equation*}
Where the two highlighted points are the original two points.
Now if we take the other two points as the points in our
mesh pattern then we get an occurrence of the pattern
we originally desired, and hence the two patterns are coincident.
It is also possible to calculate this coincidence by an extension of the Third
Dominating rule, where we allow a sequence of point addition operations, this
is discussed further in the future work section.

The other reasoning applies to the patterns
\begin{equation*}
    m_1 = \pattern{}{1,2}{0/1,0/2,2/0,2/1} \text{ and } m_2 = \pattern{}{1,2}{0/1,0/2,1/1,2/0,2/1}
\end{equation*}
which are coincident by experimentation.

In order to prove this coincidence we will proceed by mathematical
induction on the number of points in region corresponding to the middle box.
We call this number \(n\).
\begin{description}
    \item [Base Case \((n=0)\)] The base case holds since we can freely shade the
        box if it contains no points.
    \item [Inductive Hypothesis \((n=k)\)] Suppose that we can find an occurrence
        of the second pattern if we have an occurrence of the first with \(k\) points
        in the middle box.
    \item [Inductive Step \((n=k+1)\)] Suppose that we have \((k+1)\) points in
        the middle box. Choose the bottom most point in the middle box, giving the
        mesh pattern
     \begin{equation*}
        \begin{tikzpicture}[scale=0.5]
            \modpattern[5]{}{1,2,3}{0/1,0/2,0/3,1/1,2/1,3/0,3/1,3/2}
            \draw (2.5,0.5) node {\(X\)};
        \end{tikzpicture}
    \end{equation*}
        Now we need to consider the box labelled \(X\). If this box is empty then
        we have an occurrence of \(m_2\) and are done. If this box contains
        any points then we gain some extra shading on the mesh pattern
        as any points in these boxes would create an occurrence of the dominating pattern \(\perm{2,3,1}\)
     \begin{equation*}
        \begin{tikzpicture}[scale=0.5]
            \modpattern[5]{2,3}{1,2,3}{0/1,0/2,0/3,1/1,1/2,1/3,2/1,3/0,3/1,3/2}
            \filldraw (2.5,0.5) circle (4 pt);
        \end{tikzpicture}
    \end{equation*}
        The two highlighted points form an occurrence of \(m_1\) with \(k\)
        points in the middle box, and thus by the Inductive Hypothesis we
        are done.
\end{description}
By induction we have that every occurrence of \(m_1\) leads to an occurrence of \(m_2\)
and by \NoteRef{not:downcmesh} every occurrence of \(m_2\) is an occurrence of
\(m_1\) so the two patterns are coincident.
This argument applies to another two pairs of classes. Therefore in
total in \(\av{\{\perm{3,2,1}, \mperm{1,2}{R}\}}\) there are \(213\) coincidence
classes.
